\documentclass[11pt]{article}
\usepackage[T1]{fontenc}
\usepackage{lmodern}
\usepackage[margin=1in]{geometry}
\usepackage{amsmath,amssymb,amsthm,mathtools}
\usepackage{booktabs,array}
\usepackage{enumitem}
\usepackage{xcolor}
\usepackage{microtype}
\usepackage[colorlinks=true,linkcolor=blue!55!black,citecolor=blue!55!black,urlcolor=blue!55!black]{hyperref}
\usepackage[nameinlink,noabbrev]{cleveref}

\hypersetup{
  pdftitle={The General-Rank Square Endpoint for MCG-Finite Surface-Group Representations},
  pdfauthor={Robert Sneiderman},
  pdfsubject={New, unrefereed candidate note},
  pdfkeywords={mapping class groups, surface groups, local systems, endpoint estimates}
}

\newtheorem{theorem}{Theorem}[section]
\newtheorem{proposition}[theorem]{Proposition}
\newtheorem{lemma}[theorem]{Lemma}
\newtheorem{corollary}[theorem]{Corollary}
\theoremstyle{definition}
\newtheorem{remark}[theorem]{Remark}

\newcommand{\C}{\mathbf C}
\newcommand{\Ocal}{\mathcal O}
\newcommand{\Vcal}{\mathcal V}
\newcommand{\Mod}{\operatorname{Mod}}
\newcommand{\rk}{\operatorname{rk}}
\newcommand{\ad}{\operatorname{ad}}
\newcommand{\GL}{\operatorname{GL}}

\title{The General-Rank Square Endpoint for\\
MCG-Finite Surface-Group Representations}
\author{Robert Sneiderman}
\date{Version 0.1.0-candidate\\13 August 2026}

\begin{document}
\maketitle

\begin{abstract}
Landesman--Litt prove that an MCG-finite representation of a punctured
genus-$g$ surface group into $\GL_r(\C)$ has finite image when
$r<\sqrt{g+1}$.  This note gives a candidate extension to the square
endpoint $r^2=g+1$.  The new input is an equality-case refinement of their
parabolic multiplication estimate.  It extends their Artin vanishing from
coefficient rank strictly below $g$ through rank $g$.  We then audit the
resulting substitutions in their rigidity, integrality, nonabelian-Hodge,
characteristic-extension, and socle arguments.  This is a new, unrefereed
candidate proof, not an established theorem.
\end{abstract}

\begin{center}
\fbox{\parbox{0.92\textwidth}{
\textbf{Status and scope.}  The statement $g\ge r^2-1$ is
\textbf{CANDIDATE}.  The published strict range is due to
Landesman--Litt.  This note does not establish the sharper proposed range
$g\ge r^2-4$, and it makes no new genus-three or genus-four rank-three
claim.  Successful compilation and finite arithmetic checks do not replace
independent mathematical review.}}
\end{center}

\begin{center}
\small Copyright \textcopyright\ 2026 Robert Sneiderman.  Licensed under
\href{https://creativecommons.org/licenses/by/4.0/}{CC BY 4.0}.
\end{center}

\tableofcontents

\section{Statement and published inputs}\label{sec:inputs}

Let $\Sigma_{g,n}$ be an orientable genus-$g$ surface with $n$ punctures.
The mapping class group $\Mod_{g,n}$ acts by precomposition on conjugacy
classes of representations of $\pi_1(\Sigma_{g,n})$.  A representation is
\emph{MCG-finite} if its conjugacy class has finite orbit.

\begin{theorem}[Square-endpoint finite-image theorem]
\label{thm:square-endpoint}
Let $g\ge3$, $n\ge0$, and $r^2\le g+1$.  Every MCG-finite representation
\[
 \rho:\pi_1(\Sigma_{g,n})\longrightarrow\GL_r(\C)
\]
has finite image.
\end{theorem}

The strict case $r^2<g+1$ is Landesman--Litt
\cite[Theorem~1.2.1]{LLCan}.  Only $r^2=g+1$ is new here.  The proof uses the
following published statements in their cited forms.

\begin{enumerate}[label=\textup{(I\arabic*)},leftmargin=2.9em]
\item\label{input:pairing}
Landesman--Litt \cite[Proposition~5.2.3]{LLCan}: if $E_*$ is a semistable
parabolic bundle of parabolic degree zero and rank $d$, and multiplication
by a nonzero
\[
 v\in H^0(C,E\otimes\omega_C(D))
\]
has rank $s$ on $H^0(C,E^\vee\otimes\omega_C)$, then $d\ge g-s$.

\item\label{input:hn}
The proof of Landesman--Litt \cite[Proposition~6.3.6]{LLGeo}, together with
\cite[Lemma~6.2.3]{LLGeo}: for the coparabolic bundle used in
\cref{input:pairing}, a generically generated subbundle $A$ of corank $c$
and section deficit $\delta$ gives $d\ge gc-\delta$.  After quotienting the
HN part $N^t$ of slope greater than $2g-2$, the same proof gives
\[
 \rk(F/N^t)\ge gc-\delta.
\]
When $N^t=0$, its complete inequality chain is
\begin{equation}
 \deg F+(1-g)\rk F
 \le \frac{\deg A}{2}+\rk A+\delta
 \le \frac{\mu(F)\rk A}{2}+\rk A+\delta.
 \label{eq:hn-chain}
\end{equation}

\item\label{input:hodge}
Landesman--Litt \cite[Lemma~6.1.1 and Theorem~1.7.1]{LLCan}: if $U$ is an
irreducible unitary fibral local system and
$0\ne L\subset R^1\pi_*^\circ U$ is irreducible, then, after possible
complex conjugation, a nonzero Hodge vector has multiplication rank at most
$\rk L/2$; in rank one one has $\rk L\ge2g-2$.

\item\label{input:artin}
Landesman--Litt \cite[Lemma~2.4.2 and Theorem~6.2.1]{LLCan}: after a
dominant etale base change, a free Artin local system whose fibral
restriction is a constant deformation of a unitary system decomposes into
terms
\[
 U_i\otimes\pi^*W_i,
\]
with $U_i$ irreducible unitary on the fibres and $W_i$ free over the Artin
algebra.  Their fixed-part bound proves vanishing when the total coefficient
rank is strictly less than $g$.

\item\label{input:pipeline}
Landesman--Litt \cite[Sections~8.2--8.5]{LLCan}: adjoint fixed-part
vanishing gives strong cohomological rigidity and formal constancy.  The
remaining projective and linear integrality, unitary finiteness, and
nonabelian-Hodge steps use no further occurrence of
$r<\sqrt{g+1}$; the PVHS-to-unitary input is available in the larger range
$r<2\sqrt{g+1}$ \cite[Theorem~1.2.13]{LLGeo}.  The canonical-representations
paper cites the earlier numbering Theorem~1.2.12.

\item\label{input:socle}
Landesman--Litt \cite[Theorem~7.2.1, Lemma~8.6.1, and
Section~8.7]{LLCan}: for an irreducible finite-group representation
$\sigma$, every nonzero finite-index-mapping-class-stable subspace of the
relevant isotypic cohomology has dimension at least
$2g-2\dim\sigma$.  Comparing this with characteristic extension support and
then inducting on the socle filtration proves the non-semisimple case in
the published range.
\end{enumerate}

For clarity about conventions, if $E_*$ has canonical weights in $[0,1)$,
the parabolic bundle to which \cref{input:hn} is applied is
\[
 (E_*)^\vee\otimes\omega_C(D).
\]
Its coparabolic degree-zero underlying ordinary bundle is
$F=E^\vee\otimes\omega_C$.  All ordinary section counts below concern this
$F$.

\section{The strict equality-case estimate}\label{sec:endpoint}

\begin{proposition}[Strict stable pairing estimate]
\label{prop:endpoint-pairing}
Let $C$ be a smooth projective curve of genus $g\ge2$, let $D$ be reduced,
and let $E_*$ be parabolically stable of parabolic degree zero and rank
$d\ge2$.  For nonzero
\[
 v\in H^0(C,E\otimes\omega_C(D)),
\]
let $s$ be the rank of the multiplication map
\[
 H^0(C,E^\vee\otimes\omega_C)
 \longrightarrow H^0(C,\omega_C^2(D)).
\]
Then
\[
 \boxed{d\ge g-s+1.}
\]
\end{proposition}

\begin{proof}
The assertion is automatic when $s\ge g$, so suppose $s<g$.  Input
\ref{input:pairing} gives $d\ge g-s$.  Assume for contradiction that
equality holds:
\begin{equation}
 d+s=g.
 \label{eq:equality}
\end{equation}

Put $F=E^\vee\otimes\omega_C$.  Multiplication by $v$ gives a nonzero map
from $F$ to the line bundle $\omega_C^2(D)$.  Let $K$ be its kernel.  The
image is torsion-free, so $K$ is a saturated corank-one subbundle of $F$.
Let $A\subset F$ be the saturation of the evaluation image of $H^0(C,K)$.
Saturation stays in $K$ because $K$ is saturated.  Every section of $K$
factors through the evaluation image, while $A\subset K$, so
\[
 H^0(C,A)=H^0(C,K).
\]
Consequently, with $c=d-\rk A$ and
$\delta=h^0(C,F)-h^0(C,A)$,
\[
 \delta=h^0(C,F)-h^0(C,K)=s.
\]

Input \ref{input:hn} now gives
\[
 d\ge gc-s.
\]
Together with \eqref{eq:equality} and $c\ge1$, this forces $c=1$.

Let $N^t$ be the sum of the HN pieces of $A$ with slope greater than
$2g-2$.  The proof of Proposition 6.3.6 preserves the deficit because
$H^1(C,N^t)=0$, and yields
\[
 \rk(F/N^t)\ge gc-s=g-s=d.
\]
Since $F$ has rank $d$, it follows that $N^t=0$.

The intermediate inequality in the proof of Lemma 6.2.3 is therefore
\begin{equation}
 \mu(F)+2\le (2g-\mu(F))d+2s.
 \label{eq:slope}
\end{equation}
The coparabolic quotient-slope lemma gives $\mu(F)\ge2g-2$.  At
$\mu(F)=2g-2$, the two sides of \eqref{eq:slope} are equal by
\eqref{eq:equality}.  Increasing $\mu(F)$ increases the left side and
decreases the right side.  Hence
\[
 \mu(F)=2g-2.
\]

Write $S$ for the total canonical parabolic weight of $E_*$.  Since
$\deg E=-S$,
\[
 \mu(F)=2g-2+\frac{S}{d}.
\]
Thus $S=0$.  The canonical weights are nonnegative, so every weight is zero.
Because the weights in a parabolic flag are strictly increasing, each flag
is the trivial one-step flag.  Parabolic stability is now ordinary
stability of the degree-zero bundle $E$; dualizing and twisting show that
$F$ is ordinary stable of slope $2g-2$.

We may finally use the complete chain \eqref{eq:hn-chain}.  Here
$\rk A=d-1$, $\delta=s$, and $\mu(F)=2g-2$, so it reads
\[
 \deg F+(1-g)d
 \le \frac{\deg A}{2}+d-1+s
 \le \frac{(2g-2)(d-1)}2+d-1+s.
\]
The leftmost term is $d(g-1)$, and the rightmost term is
\[
 g(d-1)+s=d(g-1)
\]
by \eqref{eq:equality}.  Equality therefore holds throughout.  In
particular,
\[
 \deg A=(2g-2)(d-1),
 \qquad \mu(A)=\mu(F).
\]
But $d\ge2$ makes $A$ a nonzero proper subbundle of the stable bundle $F$,
a contradiction.
\end{proof}

\begin{remark}[The rank-one exception]
The condition $d\ge2$ is necessary.  If $D=\varnothing$ and
$E\in\operatorname{Pic}^0(C)$ is nontrivial, then
$h^0(E\omega_C)=h^0(E^{-1}\omega_C)=g-1$.  Multiplication by a nonzero
section of $E\omega_C$ is injective on $H^0(E^{-1}\omega_C)$, so
$s=g-1$ and the published weak bound is sharp.
\end{remark}

\section{Fixed parts and Artin vanishing through rank
\texorpdfstring{$g$}{g}}

\begin{corollary}[Strict fixed-part estimate]
\label{cor:fixed-part}
Let $U$ be an irreducible unitary fibral local system of rank $d$ on a
punctured versal genus-$g$ family.  Every nonzero irreducible sub-local
system $L\subset R^1\pi_*^\circ U$ satisfies
\[
 \rk L\ge
 \begin{cases}
 2g-2,&d=1,\\
 2g-2d+2,&d\ge2.
 \end{cases}
\]
\end{corollary}

\begin{proof}
For $d=1$ this is input \ref{input:hodge}.  For $d\ge2$, that input produces
a nonzero Hodge vector whose multiplication rank $s$ is at most
$\rk L/2$.  Proposition~\ref{prop:endpoint-pairing} gives
$s\ge g-d+1$, hence $\rk L\ge2g-2d+2$.
\end{proof}

\begin{corollary}[Artin vanishing through rank $g$]
\label{cor:artin-endpoint}
Let $g\ge3$.  In the setup of Landesman--Litt
\cite[Theorem~6.2.1]{LLCan}, its conclusion remains valid when
$\rk_A\Vcal\le g$.
\end{corollary}

\begin{proof}
Use input \ref{input:artin}.  It suffices to treat a summand
$U\otimes\pi^*W$, where $u=\rk U$, $w=\rk W$, and $uw\le g$.  The
maximal-ideal filtration reduces to the complex residual coefficient
$W_0$.  A nonzero invariant gives a nonzero morphism
\[
 W_0^\vee\longrightarrow R^1\pi_*^\circ U.
\]
Choose an irreducible sub-local system in its nonzero image.  Its rank is at
most $w$, so Corollary~\ref{cor:fixed-part} gives
\[
 w\ge
 \begin{cases}
 2g-2,&u=1,\\
 2g-2u+2,&u\ge2.
 \end{cases}
\]

If $w=1$, the displayed lower bound is at least two for every $u\le g$, a
contradiction.  If $w\ge2$ and $u=1$, then
$uw=w\ge2g-2>g$.  Finally, if $w\ge2$ and $u\ge2$, the constraint $uw\le g$
gives $u\le g/2$, while
\[
 uw\ge2u(g-u+1)\ge2g>g.
\]
Every case contradicts $uw\le g$.
\end{proof}

\section{The square-endpoint pipeline}\label{sec:pipeline}

We now prove Theorem~\ref{thm:square-endpoint}.  The following table records the
only substitutions in the published Landesman--Litt proof.

\begin{center}
\small
\begin{tabular}{>{\raggedright\arraybackslash}p{0.22\textwidth}
                >{\raggedright\arraybackslash}p{0.31\textwidth}
                >{\raggedright\arraybackslash}p{0.36\textwidth}}
\toprule
Published step & Strict-rank use & Endpoint replacement\\
\midrule
Proposition 8.2.1 & Theorem 6.2.1 for the rank-$r^2-1$ adjoint &
Corollary~\ref{cor:artin-endpoint}; Schur and Leray are unchanged\\
Lemma 8.3.3 & Strong rigidity from Proposition 8.2.1 &
Endpoint strong rigidity; no further rank estimate\\
Lemma 8.3.4 & Projective integrality &
Endpoint projective integrality; the finite-determinant lift is unchanged\\
Proposition 8.4.1 & Rigidity, integrality, and the larger isomonodromy range &
Endpoint rigidity and integrality; $r<2\sqrt{g+1}$ is automatic\\
Lemma 8.5.1 & Theorem 6.2.1 on every Artin truncation of the adjoint &
Corollary~\ref{cor:artin-endpoint}\\
Lemma 8.5.2 & Lemma 8.5.1 & No additional numerical hypothesis\\
Lemma 8.6.1 & Strict AM--GM bound on extension support &
The endpoint comparison in Lemma~\ref{lem:extensions}\\
Section 8.7 & Semisimple theorem and Lemma 8.6.1 &
The same socle induction\\
\bottomrule
\end{tabular}
\end{center}

\begin{proposition}[Semisimple square endpoint]
\label{prop:semisimple}
Assume $r^2=g+1$.  Every semisimple MCG-finite rank-$r$ representation has
finite image.
\end{proposition}

\begin{proof}
If the representation is reducible, each irreducible summand has rank
strictly less than $r=\sqrt{g+1}$.  Semisimple subquotients remain
MCG-finite \cite[Proposition~2.1.1]{LLCan}, so the published strict theorem
makes every summand finite.

Assume the representation is irreducible.  Its trace-free adjoint has rank
\[
 r^2-1=g.
\]
Thus Corollary~\ref{cor:artin-endpoint} replaces Theorem 6.2.1 in both strong
cohomological rigidity \cite[Proposition~8.2.1]{LLCan} and formal constancy
\cite[Lemma~8.5.1]{LLCan}.  The projective and linear integrality arguments,
unitary finiteness, and passage from formal to global constancy then apply as
in \cite[Sections~8.3--8.5]{LLCan}.  The nonabelian-Hodge deformation ends
at a unitary system because the required inequality is
\[
 r<2\sqrt{g+1}=2r.
\]
The endpoint is therefore conjugate to a finite unitary representation.
\end{proof}

\begin{lemma}[Characteristic extensions at equality]
\label{lem:extensions}
Assume $r^2=g+1$.  Let an MCG-finite representation be an extension of two
semisimple finite-image representations of positive ranks $a,b$ with
$a+b\le r$, and suppose the subrepresentation is characteristic.  Then the
extension splits and has finite image.
\end{lemma}

\begin{proof}
Decompose the finite-image coefficient as
\[
 \rho_2^\vee\otimes\rho_1
 =\bigoplus_i\sigma_i^{\oplus n_i},
 \qquad d_i=\dim\sigma_i.
\]
The AM--GM inequality gives
\[
 n_i d_i\le ab
 \le\frac{(a+b)^2}{4}
 \le\frac{r^2}{4}
 =\frac{g+1}{4}.
\]
In particular $d_i\le(g+1)/4<g$.  After passing to the finite Galois cover
that trivializes both factors, the construction in
\cite[Lemma~8.6.1]{LLCan} turns any nonzero projected extension into a
finite-index-mapping-class-stable quotient of the relevant
$\sigma_i$-isotypic homology component of dimension at most $n_i d_i$.
Dualizing embeds the dual quotient as a stable subrepresentation of the
dual $\sigma_i^\vee$-isotypic cohomology component, with the same dimension;
here $\dim\sigma_i^\vee=d_i$.  Input \ref{input:socle} therefore gives the
lower bound
\[
 2g-2d_i
 \ge 2g-\frac{g+1}{2}
 =\frac{3g-1}{2}
 >\frac{g+1}{4}
 \ge n_i d_i,
\]
a contradiction.  Hence every projected map vanishes on the finite cover.
The representation is therefore trivial on a finite-index subgroup and has
finite image.  Over $\C$, Maschke semisimplicity then splits the extension.
\end{proof}

\begin{proof}[Proof of Theorem~\ref{thm:square-endpoint}]
The strict case is published.  At equality, the semisimple case is
Proposition~\ref{prop:semisimple}.  For an arbitrary representation, apply the
characteristic socle filtration exactly as in
\cite[Section~8.7]{LLCan}.  Every proper factor has rank strictly less than
$r=\sqrt{g+1}$ and is already finite by the published theorem; equivalently,
one may induct on the length of the socle filtration.  Lemma
\ref{lem:extensions} splits each characteristic extension.  The original
representation therefore has finite image.
\end{proof}

\section{Audit boundary}

The new mathematics in this note is concentrated in three places:

\begin{enumerate}
\item the equality case in Proposition~\ref{prop:endpoint-pairing}, especially the
deduction $N^t=0$, the forced vanishing of every parabolic weight, and the
final equality-of-slopes contradiction;
\item the rank-$g$ endpoint in Corollary~\ref{cor:artin-endpoint}; and
\item the strict support margin in Lemma~\ref{lem:extensions}.
\end{enumerate}

The rest is a source-level substitution into the published pipeline.  This
note accepts the cited Hodge, coparabolic, integrality, isomonodromy, and
Putman--Wieland inputs in their published forms and does not reprove them.
The equality endpoint proved here remains a new, unrefereed candidate.

The first square pairs are
\[
 (g,r)=(3,2),(8,3),(15,4),(24,5),\ldots.
\]
The proof does not use a Spin lift and therefore includes even rank at the
square endpoint.  It supplies no mechanism below $g=r^2-1$.

\paragraph{Assistance disclosure.}
OpenAI Codex (GPT-5.6) assisted with preparation and internal checking of this
manuscript.  The author reviewed the manuscript and is responsible for all
statements, proofs, and errors.

\begin{thebibliography}{9}

\bibitem{LLCan}
A.~Landesman and D.~Litt,
\emph{Canonical representations of surface groups},
Ann. of Math. (2) \textbf{199} (2024), 823--897;
\href{https://arxiv.org/abs/2205.15352}{arXiv:2205.15352v4}.

\bibitem{LLGeo}
A.~Landesman and D.~Litt,
\emph{Geometric local systems on very general curves and isomonodromy},
J. Amer. Math. Soc. \textbf{37} (2024), 683--729;
\href{https://arxiv.org/abs/2202.00039}{arXiv:2202.00039v3}.

\end{thebibliography}

\end{document}
